\documentclass[12pt, a4paper]{article}

% --- Базовые пакеты ---
\usepackage[utf8]{inputenc}
\usepackage[russian]{babel}
\usepackage[T2A]{fontenc}
\usepackage{geometry}
\usepackage{xcolor}
\usepackage{listings}
\usepackage{longtable}
\usepackage{colortbl} % Для раскраски строк в таблицах
\usepackage{alltt}    % Для консольного вывода

\geometry{top=2cm, bottom=2cm, left=2cm, right=1.5cm}

% --- Настройка стиля кода (как на скриншоте) ---
\definecolor{codegreen}{rgb}{0,0.6,0}
\definecolor{codegray}{rgb}{0.5,0.5,0.5}

\lstset{
    language=C++,
    basicstyle=\ttfamily\small,
    keywordstyle=\color{blue},
    commentstyle=\color{codegreen},
    stringstyle=\color{red},
    % Настройки нумерации
    numbers=left,
    numberstyle=\tiny\color{black}, % Цвет номеров как в оригинале
    numbersep=5pt,                  % Расстояние от номера до кода
    % Настройки рамки (двойная линия слева)
    frame=leftline,
    framerule=1pt,                  % Толщина линии
    framesep=2pt,                   % Отступ линии от контента
    xleftmargin=15pt,               % Сдвиг всего блока вправо
    % Дополнительно
    breaklines=true,
    showstringspaces=false,
    tabsize=4
}

\newcommand{\CWHeader}[1]{\section*{#1}}
\newcommand{\CWProblem}[1]{\subsection*{Задание} #1}

\begin{document}

    \input{src/titlepage}
    \CWHeader{Лабораторная работа \textnumero 2}

\CWProblem{
Требуется разработать программную библиотеку, реализующую словарь на основе 
структуры данных PATRICIA (Practical Algorithm To Retrieve Information Coded In Alphanumeric), 
а также программу для работы с данным словарём.

Словарь хранит пары <<ключ-значение>>, где ключом является регистронезависимая 
строка, состоящая из букв английского алфавита длиной не более 256 символов, 
а значением — целое число в диапазоне от 0 до $2^{64} - 1$. Разным ключам может 
соответствовать одинаковое значение.

Программа должна обрабатывать входной поток построчно до конца файла. Поддерживаются 
следующие операции:

\begin{itemize}
    \item \texttt{+ word value} — добавить слово в словарь. \\
    Вывод: \texttt{OK}, если вставка успешна; \texttt{Exist}, если ключ уже существует.
    
    \item \texttt{- word} — удалить слово из словаря. \\
    Вывод: \texttt{OK}, если слово было удалено; \texttt{NoSuchWord}, если слово не найдено.
    
    \item \texttt{word} — поиск слова в словаре. \\
    Вывод: \texttt{OK: value}, если слово найдено; \texttt{NoSuchWord}, если отсутствует.
    
    \item \texttt{! Save path} — сохранить словарь в бинарном виде в файл. \\
    Вывод: \texttt{OK} при успехе или \texttt{ERROR: <message>} при ошибке.
    
    \item \texttt{! Load path} — загрузить словарь из файла. \\
    При успешной загрузке текущий словарь заменяется. \\
    Вывод: \texttt{OK} или \texttt{ERROR: <message>} при ошибке.
\end{itemize}

Во всех случаях возникновения системных ошибок программа должна выводить сообщение 
в формате \texttt{ERROR: <message>}.

{\bfseries Вариант структуры данных:} PATRICIA (битовое сжатое префиксное дерево).

{\bfseries Ограничения:}
\begin{itemize}
    \item Время работы: не более 15 секунд
    \item Память: не более 512 Мб
    \item Использование стандартных ассоциативных контейнеров запрещено
\end{itemize}
}
    \section{Исходный код}

Программа реализует словарь на основе структуры данных PATRICIA в битовом представлении 
(bitwise PATRICIA), соответствующем классическому описанию Horowitz и Sahni.

В данной реализации ключи рассматриваются как последовательности битов, а навигация по дереву 
осуществляется на основе значения отдельных битов ключа. Каждая вершина хранит индекс бита, 
по которому происходит ветвление.

Ключевые особенности реализации:
\begin{itemize}
    \item Внутренние вершины не хранят символы строк, а содержат индекс проверяемого бита;
    
    \item Переход к левому или правому потомку осуществляется в зависимости от значения 
    соответствующего бита ключа;
    
    \item Листовые вершины содержат полный ключ и соответствующее значение;
    
    \item Используется специальная фиктивная вершина (header), упрощающая обработку граничных случаев.
\end{itemize}

Основные операции:
\begin{itemize}
    \item \texttt{insert} — вставка нового ключа на основе первой различающейся позиции в битовом представлении;
    \item \texttt{find} — поиск ключа путём последовательного сравнения битов;
    \item \texttt{erase} — удаление ключа с перестройкой связей дерева;
    \item \texttt{save/load} — сохранение и загрузка словаря в бинарном формате.
\end{itemize}

Сложность операций:
\begin{itemize}
    \item Поиск, вставка и удаление выполняются за $O(L)$, где $L$ — длина ключа в битах;
    \item В отличие от символьных деревьев, отсутствует посимвольное сравнение строк.
\end{itemize}

\begin{lstlisting}[breaklines=true, postbreak=\mbox{\textcolor{red}{$\hookrightarrow$}\space}]
#include <iostream>
#include <vector>
#include <string>
#include <algorithm>
#include <fstream>
#include <sstream>
#include <cstdint>

using namespace std;

struct Node {
    string key;
    uint64_t value;
    int bitIndex;
    Node *left, *right;

    Node(string k, uint64_t v, int b) 
        : key(k), value(v), bitIndex(b), left(this), right(this) {}
};

class Patricia {
private:
    Node* header;

    static string toLower(string s) {
        for (char& c : s) c = (char)tolower((unsigned char)c);
        return s;
    }

    static int getBit(const string& s, int i) {
        if (i < 0) return 0;
        size_t byteIdx = i / 8;
        int bitIdx = 7 - (i % 8);
        if (byteIdx >= s.size()) return 0;
        return (s[byteIdx] >> bitIdx) & 1;
    }

    static int firstDiffBit(const string& a, const string& b) {
        int maxLen = (int)max(a.size(), b.size()) * 8;
        for (int i = 0; i < maxLen; ++i) {
            if (getBit(a, i) != getBit(b, i)) return i;
        }
        return maxLen;
    }

    Node* search(const string& key) const {
        Node* prev = header;
        Node* curr = header->left;
        while (prev->bitIndex < curr->bitIndex) {
            prev = curr;
            curr = getBit(key, curr->bitIndex) ? curr->right : curr->left;
        }
        return curr;
    }

    void clear(Node* n) {
        if (!n) return;
        if (n->left->bitIndex > n->bitIndex) clear(n->left);
        if (n->right->bitIndex > n->bitIndex) clear(n->right);
        delete n;
    }

    void collect(Node* n, vector<pair<string, uint64_t>>& data) const {
        if (!n) return;
        data.push_back({n->key, n->value});
        if (n->left->bitIndex > n->bitIndex) collect(n->left, data);
        if (n->right->bitIndex > n->bitIndex) collect(n->right, data);
    }

public:
    Patricia() {
        header = new Node("", 0, -1);
        header->left = header; 
    }

    ~Patricia() {
        if (header->left != header) clear(header->left);
        delete header;
    }

    bool find(const string& k, uint64_t& val) const {
        string key = toLower(k);
        Node* res = search(key);
        if (res != header && res->key == key) {
            val = res->value;
            return true;
        }
        return false;
    }

    bool insert(const string& k, uint64_t val) {
        string key = toLower(k);
        Node* found = search(key);
        if (found != header && found->key == key) return false;

        int diff = firstDiffBit(key, found->key);
        Node* prev = header;
        Node* curr = header->left;

        while (prev->bitIndex < curr->bitIndex && curr->bitIndex < diff) {
            prev = curr;
            curr = getBit(key, curr->bitIndex) ? curr->right : curr->left;
        }

        Node* newNode = new Node(key, val, diff);
        if (getBit(key, diff)) {
            newNode->left = curr;
            newNode->right = newNode;
        } else {
            newNode->left = newNode;
            newNode->right = curr;
        }

        if (getBit(key, prev->bitIndex)) prev->right = newNode;
        else prev->left = newNode;

        return true;
    }

    bool erase(const string& k) {
        string key = toLower(k);
        Node *gp = nullptr, *p = header, *c = header->left;

        while (p->bitIndex < c->bitIndex) {
            gp = p; p = c;
            c = getBit(key, c->bitIndex) ? c->right : c->left;
        }

        if (c == header || c->key != key) return false;

        if (p != c) {
            c->key = p->key;
            c->value = p->value;

            Node *pp = p, *cc = getBit(p->key, p->bitIndex) ? p->right : p->left;
            while (pp->bitIndex < cc->bitIndex) {
                pp = cc;
                cc = getBit(p->key, cc->bitIndex) ? cc->right : cc->left;
            }
            if (getBit(p->key, pp->bitIndex)) pp->right = c;
            else pp->left = c;
        }

        Node* target = (getBit(key, p->bitIndex) ? p->left : p->right);
        if (getBit(key, gp->bitIndex)) gp->right = target;
        else gp->left = target;

        delete p;
        return true;
    }

    bool save(const string& path) const {
        ofstream out(path, ios::binary);
        if (!out) return false;

        vector<pair<string, uint64_t>> data;
        if (header->left != header) collect(header->left, data);

        uint64_t size = data.size();
        out.write((char*)&size, sizeof(size));
        for (auto& item : data) {
            uint64_t len = item.first.size();
            out.write((char*)&len, sizeof(len));
            out.write(item.first.c_str(), len);
            out.write((char*)&item.second, sizeof(item.second));
        }
        return true;
    }

    bool load(const string& path) {
        ifstream in(path, ios::binary);
        if (!in) return false;

        if (header->left != header) clear(header->left);
        header->left = header;

        uint64_t size;
        if (!in.read((char*)&size, sizeof(size))) return false;

        for (uint64_t i = 0; i < size; ++i) {
            uint64_t len;
            if (!in.read((char*)&len, sizeof(len))) return false;
            string k(len, ' ');
            if (!in.read(&k[0], len)) return false;
            uint64_t v;
            if (!in.read((char*)&v, sizeof(v))) return false;
            insert(k, v);
        }
        return true;
    }
};

int main() {
    ios::sync_with_stdio(false);
    cin.tie(nullptr);

    Patricia tree;
    string line;

    while (getline(cin, line)) {
        if (line.empty()) continue;
        stringstream ss(line);
        string cmd;
        ss >> cmd;

        if (cmd == "+") {
            string key; uint64_t val;
            ss >> key >> val;
            cout << (tree.insert(key, val) ? "OK" : "Exist") << "\n";
        } 
        else if (cmd == "-") {
            string key; ss >> key;
            cout << (tree.erase(key) ? "OK" : "NoSuchWord") << "\n";
        } 
        else if (cmd == "!") {
            string sub, path;
            ss >> sub >> path;
            if (sub == "Save") cout << (tree.save(path) ? "OK" : "ERROR") << "\n";
            else if (sub == "Load") cout << (tree.load(path) ? "OK" : "ERROR") << "\n";
        } 
        else {
            uint64_t val;
            if (tree.find(cmd, val)) cout << "OK: " << val << "\n";
            else cout << "NoSuchWord\n";
        }
    }

    return 0;
}
\end{lstlisting}

\pagebreak
    \section{Тест производительности}

В данном разделе проводится анализ эффективности реализованной структуры данных PATRICIA 
в сравнении со стандартным ассоциативным контейнером \texttt{std::map}.

Тестирование выполнялось на случайно сгенерированных данных. Ключи представляют собой строки 
длиной до 256 символов, состоящие из букв английского алфавита.

Замерялось суммарное время выполнения операций вставки, поиска и удаления.

\begin{alltt}
--------------------------------------------------------
| Elements   | Patricia           | std::map           |
--------------------------------------------------------
|     100000 |             107 ms |             138 ms |
|     500000 |            1034 ms |            1194 ms |
|    1000000 |            1720 ms |            2715 ms |
--------------------------------------------------------
\end{alltt}

Как видно из результатов, структура PATRICIA демонстрирует более высокую производительность 
по сравнению с \texttt{std::map}, особенно на больших объёмах данных.

Это объясняется следующими причинами:

\begin{itemize}
    \item Используется побитовое сравнение ключей вместо посимвольного;
    
    \item Отсутствует зависимость от числа элементов ($\log N$), как в сбалансированных деревьях;
    
    \item Структура эффективно разделяет ключи по первому различающемуся биту.
\end{itemize}

Таким образом, bitwise PATRICIA является эффективной структурой данных для хранения 
и поиска строковых ключей.
    \section{Выводы}
{\itshape Выполнив лабораторную работу по курсу <<Дискретный анализ>>, я изучил и реализовал 
битовую версию структуры данных PATRICIA, а также закрепил навыки работы с низкоуровневыми 
представлениями данных.}

В ходе работы я пришел к следующим выводам:
\begin{itemize}
    \item \textbf{Эффективность bitwise PATRICIA:} Использование побитового сравнения позволяет 
    минимизировать количество операций сравнения и ускорить выполнение поиска, вставки и удаления.
    
    \item \textbf{Отличие от символьных деревьев:} В отличие от классических префиксных деревьев, 
    где сравнение происходит посимвольно, в PATRICIA сравнение выполняется на уровне битов, 
    что делает структуру более компактной и быстрой.
    
    \item \textbf{Работа с указателями:} Реализация потребовала аккуратного управления указателями 
    и понимания структуры дерева, особенно при реализации операции удаления.
    
    \item \textbf{Бинарное представление данных:} Работа с битами и индексами битов позволила лучше 
    понять внутреннее устройство строк и способы их эффективной обработки.
    
    \item \textbf{Практическая производительность:} Эксперименты показали, что PATRICIA превосходит 
    \texttt{std::map} по времени выполнения операций на больших объёмах данных.
\end{itemize}

Данный опыт полезен при разработке высокопроизводительных систем, работающих с текстовыми ключами, 
а также при проектировании специализированных структур данных.
    \begin{thebibliography}{99}

    \bibitem{horowitz}
    Ellis Horowitz, Sartaj Sahni. 
    {\itshape Fundamentals of Data Structures.} --- Computer Science Press, 1976.
    
    \bibitem{knuth}
    Donald E. Knuth. 
    {\itshape The Art of Computer Programming, Volume 3: Sorting and Searching.} --- Addison-Wesley, 1998.
    
    \bibitem{patricia_original}
    D. R. Morrison. 
    {\itshape PATRICIA --- Practical Algorithm To Retrieve Information Coded In Alphanumeric.} \\
    Journal of the ACM, 15(4), 1968.
    
    \bibitem{cppreference}
    {\itshape C++ reference.} \\
    URL: \texttt{https://en.cppreference.com}
    
    \end{thebibliography}
    
    
\end{document}